%% v4/sec/overview.tex — §3 Approach Overview (NEW for TODAES v4)
%%
%% Status (2026-05-16): Round-19 unlock — leader Q2 ruling: write
%% prose to ~0.5 p first, then commission Fig. 3 once we can read the
%% prose and decide which boxes / arrows deserve the figure's
%% emphasis. Figure is a placeholder until that pass.
%%
%% Target length: 0.75-1.0 page in acmsmall single-column
%% (~350-450 words + one figure). Current draft: ~370 words.
%% AutoChip~\cite{autochip} places its system-overview section here
%% as a one-figure-one-paragraph onramp; we follow that template.

\section{Approach Overview}
\label{sec:overview}

Before treating Safe-Bridge (\S\ref{sec:platform}), the benchmark
methodology (\S\ref{sec:methodology}), and the results
(\S\ref{sec:results}) in depth, this section gives a single-figure
overview of how the three components compose into one closed loop.
We adopt an AutoChip~\cite{autochip}-style onramp: one figure, one
paragraph per component, no quantitative claims---those are
deferred to \S\ref{sec:results}.

\paragraph{The closed loop in one figure.}
Figure~\ref{fig:overview} shows the three components and the arrows
between them. The \textbf{Agent} (left) is the cloud-hosted LLM
checkpoint under evaluation. The \textbf{Safe-Bridge substrate}
(centre) is the NDA-safe wrapper that scrubs every byte crossing
the trust boundary. The \textbf{spec-contract evaluator} (right) is
the SPICE-driven oracle that converts a candidate netlist into a
PASS/FAIL signal on each spec metric. The control loop turns once
per agent iteration: the agent emits a JSON tool call, Safe-Bridge
forwards it to the foundry-bound host, SKILL/SPICE executes it,
Safe-Bridge scrubs the response, and the agent observes only the
scrubbed result.

\begin{figure}[t]
  \centering
  \fbox{\parbox{0.85\linewidth}{\centering\itshape
    \vspace{1.2em}
    [Fig. 3 PLACEHOLDER --- system-overview diagram]\\[0.5em]
    Three boxes: Agent (left, cloud LLM ckpt) $\rightarrow$
    Safe-Bridge substrate (centre, four scrub layers) $\rightarrow$
    Spec-contract evaluator (right, SPICE + YAML metrics).
    Arrows: agent JSON tool calls flow right through scrub-out;
    scrubbed SPICE results flow left through scrub-in. Trust
    boundary drawn between Agent and Safe-Bridge.
    \vspace{1.2em}}}
  \caption{System overview: cloud LLM agent (left) iterating against
    a foundry-bound spec-contract evaluator (right) through the
    Safe-Bridge NDA-safe substrate (centre). Trust boundary lies
    between the agent and Safe-Bridge; PDK content never crosses
    it. Diagram to be commissioned in TikZ once prose stabilises
    (Round-19 Q2 leader ruling).}
  \label{fig:overview}
\end{figure}

\paragraph{Agent.}
Each of the eleven checkpoints in the benchmark
(\S\ref{sec:methodology}) is a frozen cloud-LLM checkpoint accessed
through its vendor's public chat-completion or messages API. No
fine-tuning is performed; only the system prompt and the JSON tool
schema are fixed across all eleven checkpoints, so that the only
controlled variable is the \texttt{-{}-llm} flag at the
\filename{run\_benchmark.py} entry point. The agent receives the
spec contract, the scrubbed simulation history, and the tool
schema; it returns either a tool call (read a node, run a
simulation, write a netlist) or a final answer.

\paragraph{Safe-Bridge substrate.}
Safe-Bridge sits between the agent and the foundry host. Four
code-anchored scrub layers (\S\ref{sec:platform}) intercept every
byte in either direction: an allow-listed SKILL-entrypoint filter
on outbound tool calls, a regex-based PDK-token redactor on
inbound SPICE output, a numeric-range sanity check on returned
metrics, and a final whole-payload audit against a foundry-leak
regex. Five hard invariants (\S\ref{sec:platform}) and a
three-axis audit methodology guarantee that no PDK token, device
model name, or foundry-internal identifier ever reaches any cloud
LLM.

\paragraph{Spec-contract evaluator.}
The evaluator (\S\ref{sec:methodology}) is a thin YAML-metrics
layer over Cadence Spectre, returning a per-metric PASS/FAIL
verdict and the underlying numeric value. The metrics are
expressed in the AnalogGym~\cite{analoggym} YAML idiom but bound
to the closed industrial PDK; the cells in the contract list
(power, phase noise, tuning range, output swing, etc.) define the
33-cell PASS threshold that gates each run.

\paragraph{Why the overview matters.}
Three claims of the rest of the paper rest on the loop in
Figure~\ref{fig:overview}: (i) the cost--quality Pareto
(\S\ref{sec:results}) is a property of the agent box and the
\texttt{-{}-llm} flag only; (ii) the NDA-safe substrate claim
(\S\ref{sec:platform}) is a property of the Safe-Bridge box and
the trust boundary it enforces; and (iii) the failure-mode taxonomy
(\S\ref{sec:failures}) is the disaggregation of where the closed
loop fails when it does. The remainder of the paper drills each
box, in order.
